\documentclass[11pt,a4paper]{article}
\usepackage[margin=26mm]{geometry}
\usepackage{amsmath,amssymb,amsthm,hyperref}
\setlength{\emergencystretch}{2em}
\newtheorem{lemma}{Lemma}
\newtheorem{proposition}[lemma]{Proposition}
\newcommand{\F}{\mathbb F}
\newcommand{\Om}{\Omega}
\newcommand{\Norm}{\operatorname{N}}
\title{Orthogonal tori and their projective orders}
\author{Human Compression project}\date{}
\begin{document}\maketitle
This working note supplies the torus interface for the orthogonal part of
the generation argument. It is not yet a reviewed component. It does not
classify maximal overgroups or estimate the centralizers of reversing
involutions. In particular it does not claim the full existence theorem.

Put $E=\F_{q^{2k}}$ and $E_0=\F_{q^k}$, and give $E$, as a vector space
over $\F_q$, the quadratic form
\[
 Q(v)=\operatorname{Tr}_{E_0/\F_q}(c v^{q^k+1}),\qquad c\in E_0^\times.
\]
It is a nondegenerate form of minus type. Here is also a direct check of
the type: every nonzero value of the norm has $q^k+1$ preimages, so the
number of zeros is
$1+(q^{k-1}-1)(q^k+1)=q^{2k-1}-(q-1)q^{k-1}$.
Nondegeneracy follows from the nondegenerate trace pairing and the polar
form of the quadratic norm. Multiplication by the norm-one group is a
cyclic group of isometries of order $q^k+1$.

\begin{lemma}\label{lem:spinor}
Suppose $q$ is odd and $z$ generates this norm-one group, with $q^k+1>2$.
Multiplication by $z$ has nontrivial spinor norm.
\end{lemma}
\begin{proof}
We use the determinant formula for the spinor norm in
\cite[Definition~3.1]{MRD}; this formula and its homomorphism property are
borrowed. Since $-1$ is not an eigenvalue of multiplication by $z$, that
formula is the square class of
$\det((1+z)/2)=\Norm_{E/\F_q}(1+z)/2^{2k}$.
Choose $w\in E$ of order $2(q^k+1)$ with $w^2=z$. Such a choice exists
because $2(q^k+1)$ divides $q^{2k}-1$. Then
$w^{q^k}=-w^{-1}$. For $u=w+w^{-1}$ we have $u\ne0$ and
$u^{q^k}=-u$, and hence $u^2\in E_0$ is a nonsquare there.
Indeed a square root in $E_0$ would have to be $u$ or $-u$.
Now
\[
 \Norm_{E/E_0}(1+z)=(1+z)(1+z^{-1})=u^2.
\]
The norm from a finite extension of finite fields sends a generator of
the multiplicative group to a generator. Consequently it sends a
nonsquare to a nonsquare. The displayed determinant has nonsquare class,
as claimed. Multiplication by $z$ has determinant one, so its powers
lying in $\Om$ are precisely its even powers.
\end{proof}

\begin{lemma}\label{lem:degree}
Let $q\ge2$, $k\ge1$, and $M>2$ divide $q^k+1$. If
$M>q^{k/3}+1$, then the multiplicative order of $q$ modulo $M$ is $2k$.
\end{lemma}
\begin{proof}
That order is $2d$ with $d\mid k$ and $k/d$ odd: it divides $2k$, but
the power $q^k$ has order two. If $d<k$, then $d\le k/3$, and
$q^d\equiv q^k\equiv-1\pmod M$. Thus $M\le q^d+1$, contrary to the
assumption.
\end{proof}

The required inequalities hold throughout the ranges below. For odd
$q\ge3$ and $k\ge3$,
$(q^k+1)/4>q^{k/3}+1$: setting $y=q^{k/3}\ge3$, this is
$y^3-4y-3>0$, which follows from its value and positive derivative at
$y=3$. For even $q$ we use $M=q^k+1$ directly, with $k\ge3$.

\begin{proposition}\label{prop:orders}
Use the following natural orthogonal isometries and let $x$ denote their
projective images.
\begin{enumerate}
\item In minus dimension $2m$, $m\ge4$, put $N=q^m+1$. Use a primitive
norm-one multiplier $g$ when $q$ is even and $g^2$ when $q$ is odd.
Its projective order is $N$ in even characteristic and
$N/\gcd(4,N)$ in odd characteristic.
\item In odd dimension $2m+1$, $q$ odd and $m\ge3$, use $g^2$ on a
negative $2m$-space and the identity on its orthogonal line.
Its order is $(q^m+1)/2$.
\item In plus dimension $2m$, $m\ge4$, put
$A=q^{m-1}+1$ and $B=q+1$. On negative spaces of dimensions $2m-2$ and
$2$, use any primitive norm-one multipliers $g,h$.
For even $q$ the projective order is $AB$ when $m$ is odd and $A$ when
$m$ is even. For odd $q$ it is
\[
 \begin{cases}
 AB/\gcd(4,q^m-1),&m\text{ odd},\\
 A/2,&m\text{ even}.
 \end{cases}
\]
\end{enumerate}
All these isometries belong to the indicated omega groups.
\end{proposition}
\begin{proof}
In odd characteristic, Lemma~\ref{lem:spinor} gives omega membership in
the first two cases and cancellation of the two nontrivial spinor norms
in the third. In even characteristic the elements have odd order, so
their images under the Dickson homomorphism to a group of order two
vanish. These homomorphisms and the determinant condition defining omega
are as in \cite[Section~3]{MRD}.

The order of $g^2$ in the first case is $N/2$. The scalar $-I$ belongs
to this cyclic subgroup if and only if $4\mid N$, which gives exactly
the additional divisor stated. In odd dimension the fixed line prevents
any power other than the identity from being scalar.

For the third case, if $m$ is odd, the Euclidean algorithm gives
$\gcd(A,B)=\gcd(2,q+1)$. For odd $q$, $A$ has exactly one factor of two.
The pair $(g,h)$ has order $AB/2$; its half-order power is $(-I,-I)$
exactly when $B$ also has exactly one factor of two, namely when
$q\equiv1\pmod4$. If $m$ is even, $B\mid A$ and $A/B$ is odd.
The pair has order $A$, and for odd $q$ its $A/2$ power is $(-I,-I)$.
There is no nontrivial scalar isometry in even characteristic. These
observations prove all the projective order formulas.
\end{proof}

\begin{proposition}\label{prop:centralizer}
In the natural projective conformal orthogonal group, the centralizers
of the elements in Proposition~\ref{prop:orders} are tori of orders
$N,N,AB$, respectively. They lie in the connected adjoint component.
There is no involutory reversing automorphism represented on this
natural module by a semilinear map with nontrivial ordinary field part.
The assertion does not address the additional triality automorphisms
in plus dimension eight.
\end{proposition}
\begin{proof}
On the large negative space let $\alpha$ be the chosen multiplier, and
write $M$ for its order after identifying scalar signs. In the minus
case $M=N/\gcd(4,N)$ for odd $q$; in the plus case it is $A/2$ for odd
$q$; in even characteristic it is the full norm-one order. In the odd
dimensional case we may use the smaller value $N/\gcd(4,N)$ too.
Lemma~\ref{lem:degree} and the inequalities following it show that
$\operatorname{ord}_M(q)=2k$, where $k=m$ or $m-1$. In particular the
large block is irreducible, and is distinguished from the other blocks
by its dimension.

If a conformal map $Y$ centralizes the projective image of the chosen
isometry $X$, then $YXY^{-1}=\lambda X$. Comparing the preserved forms
gives $\lambda^2=1$. On the large block the eigenvalues of the two
similar matrices therefore give
$\alpha^{q^j}=\lambda\alpha$ for some $0\le j<2k$.
Modulo scalar signs this says $q^j\equiv1\pmod M$, so $j=0$ and
$\lambda=1$. Hence $Y$ centralizes $X$ as a linear map.

Its restriction to the large block belongs to the field algebra
$E=\F_q[X]$, and the conformal condition on multiplication by $z$ is
$\Norm_{E/E_0}(z)\in\F_q^\times$. This follows by substituting into
$Q$ and using the nondegenerate trace pairing. There are
$(q-1)(q^k+1)$ such multipliers. In minus type division by the $q-1$
scalar multipliers gives $N$.
In plus type the two block multipliers must have the same norm in
$\F_q^\times$. The number of pairs is $(q-1)AB$, so their projective
group has order $AB$. The small block is irreducible because its
primitive multiplier has order $q+1>2$; this includes $q=2$.
In odd dimension the line multiplier $a$ forces the large-block norm
to be $a^2$. Dividing both multipliers by $a$ leaves precisely the
norm-one group, of order $N$. These multiplicative centralizers are
connected tori over the algebraic closure and consequently belong to
the connected adjoint component.

Finally let $q=p^f$. An involutory automorphism with a nontrivial field
projection has $f$ even and projection $p^{f/2}$. If it reverses $x$,
comparison of the eigenvalues on the large block, again allowing the
scalar sign, implies
\[
 p^{f/2+fj}\equiv-1\pmod M
\]
for some $j$ (the sign introduced by a scalar is already quotiented
out in $M$). Squaring gives $q^{2j+1}\equiv1\pmod M$.
This contradicts $\operatorname{ord}_M(q)=2k$. Natural graph
automorphisms are represented by linear conformal maps and do not alter
this argument; exceptional triality still requires a separate argument.
\end{proof}

The distinction between a cyclic subgroup and its centralizer matters.
For example in plus type with even $m$, $x$ has order $A/2$ for odd $q$
but its adjoint torus has order $AB$. In minus dimension ten over
$\F_3$, the projective Singer image has order $61$, whereas the natural
omega Singer subgroup has order $122$. The latter convention agrees
with \cite[Table~1]{BLW}. We retain these scalar losses explicitly in
the later counting argument.

\begin{thebibliography}{9}\small
\bibitem{MRD} S.~H. Murray and C.~M. Roney-Dougal,
\emph{The spinor norm and homomorphism algorithms for classical groups},
preprint dated 2 May 2008, Definition~3.1 and Lemma~3.2.
\bibitem{BLW} D. Bubboloni, M.~S. Lucido and T. Weigel,
\emph{2-Coverings of classical groups}, arXiv:1102.0660 (2011).
\end{thebibliography}
\end{document}
